\chapter{Anwendungs-Aufgaben}\label{ch:K08_Anwendung}
\section{Kalorienverbrauch}

Folgende Beispiele sollen die meisten der bisher gelernten Programmier-Konzepte konkret veranschaulichen. Stellen Sie sich vor, Sie sind SportlerIn und möchten Ihren Kalorienbedarf berechnen, um sich für auf einen kommenden Wettkampf optimal zu ernähren. Hierzu wollen Sie zunächst Ihren theoretischen täglichen Kalorienbedarf berechnen. Der theoretische tägliche Kalorienbedarf berechnet sich wie folgt:
\begin{enumerate}
	\item Ihr \textbf{Grundumsatz} (englisch \gls{BMR}): Dies ist ihr Grundbedarf, also die Anzahl Kalorien, welche Sie theoretisch benötigen, falls Sie sich gar nicht bewegen.
	\item Ihr \textbf{Leistungsumsatz}, d.h., zusätzliche Energie, die Sie bei Aktivitäten wie Spazieren, Radfahren, Joggen usw. verbrennen.
\end{enumerate}

Beide Komponenten sind schwierig abzuschätzen. Sie möchten jedoch eine erste Einschätzung Ihrer \gls{BMR} haben, indem sie einige bekannte Formeln auf sich selber anwenden und diese vergleichen.

\begin{myexercise}\label{ex:bmr-simple}
	Schreiben Sie eine Python-Funktion, um den Grundumsatz anhand der \textbf{Harris-Benedict-Formel} zu berechnen und geben Sie das Resultat per \lstinline|return| zurück.
	\begin{enumerate}
		\item Für Männer:
		      \begin{align*}
			      \text{BMR} = & 88.362 +                                   \\
			                   & (13.397 \times \text{Gewicht in kg}) +     \\
			                   & (4.799 \times \text{Körpergrösse in cm}) - \\
			                   & (5.677 \times \text{Alter in Jahren})
		      \end{align*}
		\item Für Frauen:
		      \begin{align*}
			      \text{BMR} = & 447.593  +                                 \\
			                   & (9.247 \times \text{Gewicht in kg}) +      \\
			                   & (3.098 \times \text{Körpergrösse in cm}) - \\
			                   & (4.330 \times \text{Alter in Jahren})
		      \end{align*}
	\end{enumerate}
	\begin{myanswer}
		\lstinputlisting[language=Python]{Grundlagen_Info/00_Programmieren/Skript/Code/K08_Anwendungen/bmr_functions/grundumsatz.py}
	\end{myanswer}
\end{myexercise}

Wir haben nun eine erste, grobe Einschätzung des \gls{BMR}. Allerdings gibt es auch noch weitere mögliche Formeln, mithilfe denen der \gls{BMR} berechnet werden kann. Diese möchten wir nun ebenfalls berechnen, um die unterschiedlichen Schätzungen des \gls{BMR} anhand der verschiedenen Methoden zu vergleichen.

\begin{myexercise}\label{ex:bmr}
	Schreiben Sie zwei weitere Funktionen, um den \gls{BMR} anhand der \textbf{Mifflin-St-Jeor-Gleichung} sowie der \textbf{Katch-McArdle-Formel} zurückzugeben und berechnen Sie den Durchschnitt aller drei Formeln in einer weiteren Funktion.
	\begin{enumerate}
		\item \textbf{Mifflin-St-Jeor-Gleichung}:
		      \begin{enumerate}
			      \item Männer:
			            \begin{align*}
				            \text{BMR} & = 10 \times \text{Gewicht (kg)}        \\
				                       & \quad + 6.25 \times \text{Grösse (cm)} \\
				                       & \quad - 5 \times \text{Alter (Jahre)}  \\
				                       & \quad + 5
			            \end{align*}
			      \item Frauen:
			            \begin{align*}
				            \text{BMR} & = 10 \times \text{Gewicht (kg)}        \\
				                       & \quad + 6.25 \times \text{Grösse (cm)} \\
				                       & \quad - 5 \times \text{Alter (Jahre)}  \\
				                       & \quad - 161
			            \end{align*}
		      \end{enumerate}
		\item \textbf{Katch-McArdle-Formel} (gleich für Männer sowie Frauen):
		      
			            \begin{align*}
				            \text{BMR} & = 370                                  \\
				                       & \quad + (21.6 \times \text{FFM in kg})
			            \end{align*}
						Wobei die fettfreie Masse (FFM) zuerst berechnet werden muss:
			            \begin{align*}
				            \text{FFM} & = \text{Körpergewicht (kg)}                \\
				                       & \quad \times (1 - \text{Körperfettanteil})
			            \end{align*}

	\end{enumerate}

	\begin{myanswer}
		\lstinputlisting{Grundlagen_Info/00_Programmieren/Skript/Code/K08_Anwendungen/bmr_functions/bmr.py}
	\end{myanswer}
\end{myexercise}

Nun haben wir eine etwas verlässlichere Einschätzung unseres Grundumsatzes. Zum Grundumsatz können wir nun auch noch den Leistungsumsatz hinzufügen, um den gesamten theoretischen täglichen Energieumsatz in Kalorien zu berechnen.

\begin{myexercise}\label{ex:leistungsumsatz}
	Berechnen Sie Ihren Leistungsumsatz (durch Leistung benötigte Energie) anhand der Tabelle \ref{tab:kcal}.

	\begin{table}[H]
		\centering
		\begin{tabular}{c|*{10}{c}}
			\textbf{Gewicht (kg)}    & 60   & 65   & 70   & 75   & 80   & 85   & 90   & 95   & 100  & 105  \\
			\midrule\midrule
			\textbf{Spazieren}       & 4.5  & 5.1  & 5.4  & 5.7  & 6.0  & 6.4  & 6.7  & 7.3  & 8.0  & 9.0  \\
			\midrule
			\textbf{Schnelles Gehen} & 5.5  & 6.3  & 7.0  & 7.5  & 8.0  & 8.6  & 9.5  & 10.0 & 10.6 & 11.3 \\
			\midrule
			\textbf{Fahrradfahren}   & 6.5  & 7.5  & 8.0  & 9.0  & 9.5  & 10.3 & 11.0 & 11.7 & 12.5 & 13.6 \\
			\midrule
			\textbf{Schwimmen}       & 8.0  & 9.2  & 10.0 & 10.8 & 11.5 & 12.5 & 13.7 & 14.4 & 15.4 & 16.5 \\
			\midrule
			\textbf{Rudern}          & 11.0 & 13.0 & 14.3 & 15.0 & 16.5 & 17.5 & 19.0 & 20.5 & 21.8 & 23.5 \\
			\midrule
			\textbf{Jogging}         & 13.0 & 14.5 & 16.0 & 17.5 & 19.0 & 20.0 & 22.0 & 23.5 & 24.5 & 26.5 \\
			\midrule
			\textbf{Laufen}          & 16.0 & 18.5 & 20.0 & 22.0 & 23.5 & 25.0 & 27.5 & 29.0 & 31.0 & 33.5 \\
			\midrule
			\textbf{Sprint}          & 19.0 & 21.5 & 24.0 & 26.0 & 28.0 & 30.0 & 32.5 & 34.7 & 35.6 & 39.0 \\
		\end{tabular}
		\caption{Kalorienverbrauch für unterschiedliche Aktivitäten, pro Minute, in Abhängigkeit des Körpergewichts}
		\label{tab:kcal}
	\end{table}
	Wenn z.B. eine 80 kg schwere Person 1 Stunde lang rudert und 20 Minuten spaziert wären die Listen wie folgt:
	\begin{lstlisting}
		l_zeit = [60, 20] # Zeit in Minuten
		l_kcal = [16.5, 6] # Kalorien pro Aktivität
	\end{lstlisting}
	Erstellen Sie zwei Listen:
	\begin{enumerate}
		\item Eine Liste für die Zeit, während der Sie täglich einer Aktivität nachgehen
		\item Eine Liste mit dem Kalorienverbrauch pro Minute für jede dieser Aktivitäten
	\end{enumerate}

	Schreiben Sie eine Python-Funktion \lstinline|berechne_leistungsumsatz|, um Ihren Leistungsumsatz zu berechnen. Das erwartete Resultat für das Beispiel hier wäre: \lstinline|1110| (= Summe von \lstinline|[990, 120]|). Das Resultat soll an das Hauptprogramm zurückgegeben und in einer Variable \texttt{kalorienverbrauch} gespeichert werden.
	\begin{myanswer}
		\lstinputlisting{Grundlagen_Info/00_Programmieren/Skript/Code/K08_Anwendungen/bmr_functions/leistungsumsatz.py}
	\end{myanswer}

\end{myexercise}

\begin{myexercise}
	Schreiben Sie nun eine weitere Funktion, die Ihren gesamten Energieumsatz berechnet, als Summe der folgender Teile:
	\begin{itemize}
		\item \textbf{\gls{BMR}}, den Sie in \cref{ex:bmr} berechnet haben (Durchschnitt der drei Formeln)
		\item \textbf{Leistungsumsatz}, den Sie in \cref{ex:leistungsumsatz} berechnet haben
	\end{itemize}
	Diese Funktion ist extrem kurz und umfässt nur eine einzige (neue) Zeile.
	\begin{myanswer}
		Angenommen, die bisher geschriebenen Funktionen wurden in den Dateien \texttt{bmr.py} und \texttt{leistungsumsatz.py} gespeichert, so können Sie die Funktion \lstinline|def total(bmr, leistungsumsatz)| wie folgt schreiben:
		\lstinputlisting{Grundlagen_Info/00_Programmieren/Skript/Code/K08_Anwendungen/bmr_functions/total.py}
	\end{myanswer}
\end{myexercise}




\begin{mycaution}
	Die Zahlen zu Energiebedarf und Kalorien, die Sie im Rahmen dieser Übungen berechnet haben, sollten Sie mit Vorsicht geniessen, aus folgenden Gründen:
	\begin{itemize}
		\item Ihr realer Energiebedarf kann signifikant von Ihrem berechneten Energiebedarf abweichen. Viele weitere Faktoren, die nicht in den Formeln enthalten sind, können diesen beeinflussen, beispielsweise Ihre Köpertemperatur, Ihre Fitness, Ihre Muskelmasse, sowie Ihre \gls{NEAT}, wobei Letzteres Energie bezeichnet, die Sie durch Bewegungen, die nicht Sport sind, verbrauchen (beispielsweise durch ``Zappeln'' aber auch kurze Spaziergänge und alltägliche Aktivitäten).
		\item Sie sollten sich bezüglich konsumierter Kalorien vor Augen halten, dass ``zu viele'' konsumierte Kalorien nicht automatisch mit einer Gewichtszunahme verbunden sind. In einem Kilo Körperfett stecken ca. 8000 Kilokalorien, wobei der Köper über verschiedene Mechanismen verfügt, um ein stabiles Körpergewicht zu halten, etwa die Erhöhung oder Senkung der Körpertemperatur im Falle eines kurzfristigen Kalorien-Überschusses, respektive -Defizits. Allerdings kann es interessant sein, zu wissen, dass einige, nicht sehr sättigende Lebensmittel besonders viele Kalorien enthalten ––- wie beispielsweise hochverarbeitete Snacks oder Süssgetränke.
	\end{itemize}
\end{mycaution}
